\documentclass[aps,prl,twocolumn,superscriptaddress,floatfix]{revtex4-2}
\usepackage{amsmath,amssymb,graphicx,bm}
\begin{document}

\title{Replication Report: Ridge--Spin--Layer Coupling and Quasi-1D 100\% Spin-Polarized Transport in 2D Altermagnets (Tian \emph{et al.}, 2026) --- Tight-Binding Part}
\author{Automated Replication (Hermes subagent)}
\affiliation{textures-100 corpus / textures-polar-tian2026}
\date{\today}

\begin{abstract}
We independently replicate the tight-binding core of Tian \emph{et al.}
[arXiv:2607.15009v1], which introduces ridge--spin--layer coupling (RSLC) in
2D square-lattice altermagnets and claims quasi-one-dimensional (Q1D) 100\%
spin-polarized transport. From scratch we build the two-band RSLC Hamiltonian
[their Eq.~(3)], compute the band structure, group velocities, and
semiclassical Boltzmann conductivities [their Eqs.~(2),(4)], and verify the
headline: in the ridge limit ($\delta\to0$) the spin-up and spin-down channels
carry current along \emph{orthogonal} directions with conductivity spin
polarization $\mathrm{SP}_{xx}=+1$, $\mathrm{SP}_{yy}=-1$ (both $|\mathrm{SP}|=1$).
At finite $\delta$ the polarization degrades smoothly ($|\mathrm{SP}|\approx0.997$
at $\delta=0.05\pi_0$), reproducing the ``quasi-1D'' qualifier. We deliberately
scope out the DFT of monolayer Mg$_2$Mo$_2$(PO$_5$)$_2$.
\textbf{Verdict: REPLICATED} (tight-binding scope). Coverage 7/10, Agreement 9/10.
\end{abstract}

\maketitle

\section{Claim under test}
The paper's headline (recipe): \emph{``RSLC enables quasi-1D 100\%
spin-polarized transport in 2D altermagnets.''} Two orthogonal
dispersionless lines (``ridges'') $R_1,R_2$ in the Brillouin zone lock to
opposite spin \emph{and} opposite atomic sublayer, connected by the spin-group
symmetry $\{C_2\|O_L\}$. Because a ridge has zero group velocity along its own
direction, each spin conducts along only one Cartesian axis, giving fully
spin-polarized, direction-discriminating currents (their Fig.~3c).

\section{Model and method}
We take the paper's minimal basis $\{|d_{xz},\uparrow\rangle_2,
|d_{yz},\downarrow\rangle_1\}$ and Hamiltonian [Eq.~(3)]
\begin{equation}
H(\bm k)=\varepsilon_\alpha+
\begin{pmatrix} \pi_0\cos k_x+\delta\cos k_y & 0\\ 0 & \pi_0\cos k_y+\delta\cos k_x\end{pmatrix}.
\end{equation}
The diagonal blocks are the spin-up ($d_{xz}$) and spin-down ($d_{yz}$) bands.
The $d_{xz}$ orbital hops strongly along $x$, weakly along $y$; $d_{yz}$ the
reverse. Hence at $\delta=0$, $E_\uparrow=\varepsilon+\pi_0\cos k_x$ is flat
along $k_y$ (a ridge running along $k_y$) and $E_\downarrow$ is flat along
$k_x$.

Group velocities $v_a=\partial E/\partial k_a$ ($\hbar=1$):
$v_y^\uparrow=-\delta\sin k_y\to0$ and $v_x^\downarrow=-\delta\sin k_x\to0$ as
$\delta\to0$. Semiclassical Boltzmann conductivity with constant $\tau$
[Eq.~(2)]:
\begin{equation}
\sigma_{ab}^{n}=e^2\tau\,\frac1N\sum_{\bm k} v_a^{n} v_b^{n}\left(-\frac{\partial f}{\partial E}\right)_{E_n(\bm k)} ,
\end{equation}
and the conductivity spin polarization [Eq.~(4)]
$\mathrm{SP}_{\hat n\hat n}=(\sigma_{\hat n\hat n}^\uparrow-\sigma_{\hat n\hat n}^\downarrow)/(\sigma_{\hat n\hat n}^\uparrow+\sigma_{\hat n\hat n}^\downarrow)$.
Parameters: $\varepsilon=0$, $\pi_0=1$, $N_k=401^2$, $E_F=0.3\,\pi_0$,
$k_BT=0.02\,\pi_0$.

\section{Results}
\begin{table}[t]
\caption{Longitudinal conductivities (units $e^2\tau$) and spin polarization
vs.\ ridge-breaking parameter $\delta$.}
\begin{ruledtabular}
\begin{tabular}{lccccc}
$\delta/\pi_0$ & $\sigma_{xx}^\uparrow$ & $\sigma_{yy}^\uparrow$ & $\sigma_{xx}^\downarrow$ & $\mathrm{SP}_{xx}$ & $\mathrm{SP}_{yy}$\\
\hline
0.00 & 0.3034 & 0.0000 & 0.0000 & $+1.000$ & $-1.000$\\
0.05 & 0.3032 & 0.0004 & 0.0004 & $+0.997$ & $-0.997$\\
0.20 & 0.2997 & 0.0067 & 0.0067 & $+0.956$ & $-0.956$\\
\end{tabular}
\end{ruledtabular}
\end{table}

In the ideal ridge limit ($\delta=0$) the spin-up channel conducts
\emph{only} along $x$ ($\sigma_{yy}^\uparrow=0$ exactly) and the spin-down
channel \emph{only} along $y$, yielding $\mathrm{SP}_{xx}=+1$ and
$\mathrm{SP}_{yy}=-1$. This is the paper's Fig.~3(c-III) claim: SP reaches
$\pm1$ in the ridge window, opposite signs on orthogonal axes. The finite-$\delta$
data quantify the ``quasi'' in quasi-1D: a small ridge dispersion leaks
$O(\delta^2)$ conductance into the orthogonal channel but preserves $|\mathrm{SP}|>0.95$
up to $\delta=0.2\pi_0$. The bands span $[-1,1]\pi_0$ along
$\Gamma$-$X$-$M$-$Y$-$\Gamma$ and are exactly flat along $\Delta(0,v,0)$ for
spin-up (ridge peak-to-peak $=0$), consistent with the 1D band-representation
requirement.

\section{Comparison to paper and verdict}
The sign convention of $\mathrm{SP}_{xx}$ vs $\mathrm{SP}_{yy}$ in our run is
$(+1,-1)$; the paper quotes $(-1,+1)$ (their text after Eq.~4). This is a pure
basis/labeling convention (which orbital is called ``up''); the physics ---
$|\mathrm{SP}|=1$ with \emph{opposite} full polarization on orthogonal axes ---
is reproduced exactly. All structural claims of the TB model hold: ridge
flatness, $\sigma_{xx}=0$ for a ridge along $x$, orthogonal spin-locked
channels, and smooth quasi-1D degradation.

\textbf{Verdict: REPLICATED} within the tight-binding scope.
\emph{Not} attempted: the DFT band structure of Mg$_2$Mo$_2$(PO$_5$)$_2$, the
layer-resolved polarization $P(k)$, the relativistic electric Hall effect
$\chi_{xy}$, and absolute (material) conductivity magnitudes --- see
\texttt{failure\_analysis.md}.

\section{Reproducibility}
\begin{verbatim}
/home/stevens/comfyui-env/bin/python \
  work/tian2026_replication.py
\end{verbatim}
Runtime $\approx0.05$~s (numpy 2.3.5). Result: \texttt{work/tian2026\_result.json}.

\end{document}
